

Donc @@M0@@ est le coefficient de @@M1@@ dans la série :
@@M2@@

\textbf{Proposition 5.7.} \uline{Soient} @@M3@@ \uline{un anneau binomial}, @@M4@@ \uline{une} @@M5@@-@@M6@@-\uline{algèbre augmentée}, @@M7@@ \uline{un élément de} @@M8@@ \uline{tel que} @@M9@@ \uline{pour} @@M10@@. \uline{Alors pour tout} @@M11@@ \uline{et tout} @@M12@@, \uline{on a} :
@@M13@@

Soit @@M14@@ la @@M15@@-@@M16@@-algèbre engendrée par des générateurs @@M17@@, soumis aux relations @@M18@@ pour @@M19@@. On définit l'homomorphisme d'augmentation @@M20@@ par @@M21@@. On sait que @@M22@@ est l'anneau des polynômes à coefficients dans @@M23@@ par rapport aux indéterminées @@M24@@, avec @@M25@@, et @@M26@@, avec @@M27@@, et que la @@M28@@-filtration sur @@M29@@ est identique à la filtration par le poids, en affectant le poids @@M30@@ à @@M31@@, @@M32@@ (4.12 et 4.15).

Posons @@M33@@, @@M34@@. Nous allons d'abord montrer 5.7 pour @@M35@@ et @@M36@@. D'après ce qui précède, (5.7.1) est équivalent à :
@@M37@@
Or @@M38@@~; cette relation s'écrit donc :
@@M39@@
d'après (5.5.4). Cette relation est bien vérifiée, puisque @@M40@@ est d'augmentation nulle.

D'après 4.7, il existe un @@M41@@-@@M42@@-homomorphisme @@M43@@ compatible aux augmentations envoyant @@M44@@ sur @@M45@@, et @@M46@@ sur @@M47@@, donc @@M48@@ sur @@M49@@ et @@M50@@ 
